\chapter{Pterygium}
\index{Pterygium}
\label{sec:Pterygium}

\section{Overview}
Pterygium is a fibrovascular overgrowth of the bulbar conjunctiva that extends onto the cornea, typically originating from the nasal (medial) side. It is a common ocular surface condition associated with chronic ultraviolet light exposure, with higher prevalence in sunny climates and outdoor workers. While often benign, pterygium can cause visual disturbance when it approaches the visual axis.
\section{Classification}

\begin{description}[font=\normalfont\bfseries,leftmargin=3em,labelwidth=2.5em]
  \item[Cause] Environmental / Degenerative
  \item[Organ System] Eyes
  \item[Pathophysiology] UV-induced elastotic degeneration and fibrovascular proliferation of the conjunctiva
  \item[Duration] Chronic (months to years)
  \item[Onset] Gradual
  \item[Mode of Transmission] Non-communicable
  \item[Clinical Course] Chronic, slowly progressive
  \item[Severity] Mild to Moderate
  \item[Extent of Disease] Localized
  \item[Age Group] Adults (20--50 years)
  \item[Epidemiology] Prevalence 10--20\% in equatorial regions; higher in men
  \item[ICD-11 Code] 9A62.0
\end{description}

\section{Causes}
Chronic exposure to ultraviolet B radiation from sunlight is the primary causative factor, damaging limbal stem cells and triggering abnormal conjunctival proliferation. Genetic predisposition, chronic dry eye, and exposure to irritants like dust and wind are contributing factors. The condition is most prevalent in latitudes near the equator.

\section{Risk Factors}

\begin{itemize}[noitemsep]
  \item Chronic sunlight exposure (outdoor occupations, residence near equator)
  \item Lack of eye protection (sunglasses, hats)
  \item Male sex
  \item Older age
  \item Chronic dry eye
  \item Exposure to dust, wind, and irritants
  \item Family history
\end{itemize}

\section{Signs and Symptoms}

\subsection{Early symptoms}
\begin{itemize}[noitemsep]
  \item Early manifestations of pterygium may be subtle and nonspecific
\end{itemize}

\subsection{Common symptoms}
\begin{itemize}[noitemsep]
  \item \subsection{Common symptoms}
  \item \textbf{Common symptoms:}
  \item Fleshy, triangular growth on the conjunctiva extending onto the cornea
  \item Typically originates from the nasal side (bilateral in 10\% of cases)
  \item Redness, irritation, and foreign body sensation
  \item Tearing and dry eye symptoms
  \item Blurred vision when growth approaches the visual axis
  \item Astigmatism induced by corneal traction
  \item Pain or discomfort in the affected area
  \item Difficulty performing daily activities
\end{itemize}

\subsection{Less common symptoms}
\begin{itemize}[noitemsep]
  \item Atypical or less common presentations of pterygium may occur
\end{itemize}

\subsection{Emergency warning signs}
\begin{itemize}[noitemsep]
  \item Severe or worsening symptoms of pterygium require immediate medical evaluation
\end{itemize}
\section{Progression and Stages}
Pterygium progresses slowly over years to decades. Early lesions remain confined to the conjunctiva. As the fibrovascular tissue advances across the cornea, it can induce irregular astigmatism. Advanced pterygium may extend to or beyond the pupillary margin, significantly impairing vision. Recurrence after surgical excision is common without adjunctive therapy.

\section{Complications}

\begin{itemize}[noitemsep]
  \item Visual impairment from corneal involvement or induced astigmatism
  \item Restriction of eye movement in advanced cases
  \item Recurrence after surgical removal (up to 40--50\% without adjunctive treatment)
  \item Corneal scarring and irregularity
  \item Reduced quality of life
  \item Emotional and psychological effects
\end{itemize}

\section{Diagnosis}

\begin{itemize}[noitemsep]
  \item Clinical diagnosis by slit-lamp examination
  \item Measurement of horizontal extent onto cornea
  \item Assessment for induced astigmatism by keratometry or topography
  \item Photodocumentation for progression monitoring
  \item Lab tests to check for underlying conditions
\end{itemize}

\section{Prevention}

\begin{itemize}[noitemsep]
  \item Wear UV-blocking sunglasses and wide-brimmed hats outdoors
  \item Use artificial tears in dry or dusty environments
  \item Avoid chronic sunlight exposure during peak hours
  \item Regular ophthalmic examination for outdoor workers
  \item Go for regular check-ups so any problems are caught early
\end{itemize}

\section{Treatment Options}

\begin{itemize}[noitemsep]
  \item Observation for small, asymptomatic pterygia
  \item Lubricating eye drops for irritation and dry eye
  \item Topical corticosteroids for active inflammation
  \item Surgical excision with conjunctival autograft for vision-threatening pterygia
  \item Adjunctive mitomycin C or beta radiation to reduce recurrence
  \item Lifestyle changes and self-care
  \item Surgery when needed
\end{itemize}

\section{Living with It}

\begin{itemize}[noitemsep]
  \item Follow your treatment plan as prescribed
  \item Keep up with regular appointments with your healthcare provider
  \item Adjust your daily habits as needed (eye protection outdoors)
  \item Reach out to family, friends, or support groups for help
  \item Keep learning about your condition and how to manage it
\end{itemize}

\section{Outlook}
Pterygium is a benign condition with an excellent prognosis for vision when managed appropriately. Small lesions can be observed indefinitely. Surgical removal is curative but carries a risk of recurrence, which is significantly reduced with conjunctival autograft techniques. Long-term sun protection is essential to prevent progression and recurrence.
